This pilot study evaluates whether depth signal generalizes beyond Wikipedia by constructing generation chains from three additional domains: News articles, Academic abstracts, and Code documentation.
For each domain, we sample 200 seed texts and generate chains of depth 3 using three generators (Pythia-1B, Qwen3-1.7B, Mistral-7B).
The same 7-feature ensemble and LOGO protocol are applied.

\noindent\textbf{Important caveat:} Cross-domain generation used $t{=}1.0$, top-$p{=}1.0$, which falls outside the method's reliable operating range ($t \leq 0.7$; Appendix~\ref{app:sampling-sensitivity}). Domain and temperature effects are therefore confounded in the results below; the cross-domain signal may partly reflect temperature-induced distributional divergence rather than domain generalization alone.

\begin{table}[ht]
\centering
\small
\caption{Cross-domain depth classification under LOGO evaluation (binary: gen-1 vs.\ gen-2+). Three generators $\times$ 200 seeds per domain. Bootstrap 95\% CIs ($n$\,=\,1000).}
\label{tab:domain}
\begin{tabular}{lccc}
\toprule
\textbf{Domain} & \textbf{Binary F1} & \textbf{95\% CI} & \textbf{3-class F1} \\
\midrule
Code           & 0.754 & [0.737, 0.770] & 0.634 \\
News           & 0.746 & [0.726, 0.764] & 0.652 \\
Academic       & 0.739 & [0.720, 0.757] & 0.622 \\
\bottomrule
\end{tabular}
\end{table}

All three non-Wikipedia domains achieve binary F1 $>$ 0.73 under leave-one-generator-out evaluation. Depth signal is not an artifact of Wikipedia text properties.
The Code-domain 3-class F1 (0.634) likely reflects domain mismatch in the Wikipedia-trained reference model (Pythia-410M), whose NLL estimates may be less discriminative for code syntax; these cross-domain results are preliminary.

\paragraph{Domain $\times$ temperature interaction.}
Cross-domain generation used $t{=}1.0$, top-$p{=}1.0$ (vs.\ the main experiment's $t{=}0.7$, top-$p{=}0.9$).
Supplementary evaluation at $t{=}0.7$ reveals a domain $\times$ temperature interaction: the main Wikipedia experiment favors lower temperature (depth signal degrades monotonically with rising temperature; Appendix~\ref{app:sampling-sensitivity}), while cross-domain evaluation favors higher temperature ($t{=}1.0$ yields +3 to +11\,pp over $t{=}0.7$; News 0.746 vs.\ 0.712, Academic 0.739 vs.\ 0.628, Code 0.754 vs.\ 0.687).
This asymmetry likely reflects reference-model calibration: Pythia-410M, trained predominantly on web text, detects subtle per-step drift in Wikipedia but requires larger distributional divergence (produced by higher temperature) to detect depth signal in out-of-distribution domains.
This finding suggests that reference model selection is a key factor for cross-domain deployment of depth profiling.

\paragraph{3-class inversion as DDR evidence.}
In the cross-domain setting, 3-class F1 (0.62--0.65) falls below binary F1 (0.74--0.75), reversing the pattern observed in the main experiment (3-class 0.762 $>$ binary 0.694).
This inversion reflects degraded gen-0 vs.\ gen-1 separation: the Pythia-410M reference model is calibrated to Wikipedia text, so its baseline entropy estimates are less discriminative for out-of-domain human text and compresses the feature distance between gen-0 and gen-1.
The binary task (gen-1 vs.\ gen-2+) is unaffected because both classes are generated under identical domain and decoding conditions; recursive drift is a relative change that does not depend on absolute calibration.
This pattern provides additional evidence that DDR captures base-model generation dynamics rather than domain-specific text properties: depth signal remains robust even when the human-vs-machine boundary becomes less distinct.
